%% Drafted from the mapped corpus by scripts/draft_topics.py.
% Model: gpt-5.6-terra; source characters: 18188.
\section{Определен интеграл и суми на Дарбу}

Нека $f\colon [a,b]\to\mathbb{R}$, където $a<b$. Целта на определения интеграл е да се постави на строга основа идеята за натрупана величина: например лице под графика, изминат път при зададена скорост или общ принос на непрекъснато изменяща се функция.

\subsection{Разбиване на интервал}

\begin{defn}
\label{def:partition}
\emph{Разбиване} на интервала $[a,b]$ наричаме всяка крайна система от точки
\[
\tau\colon\quad a=x_0<x_1<\cdots<x_n=b.
\]
Точките $x_0,x_1,\ldots,x_n$ се наричат \emph{делящи точки} на разбиването. С разбиването са свързани подинтервалите
\[
[x_{i-1},x_i],\qquad i=1,\ldots,n,
\]
чийто съюз е $[a,b]$ и два от които или не се пресичат, или имат обща само крайна точка.
\end{defn}

\begin{defn}
\emph{Диаметър} на разбиването $\tau$ се нарича числото
\[
d(\tau)=\max_{1\leq i\leq n}(x_i-x_{i-1}).
\]
\end{defn}

Колкото по-малък е диаметърът, толкова по-фино е разбит интервалът. При изследване на интегруемостта ще използваме разбивания, за които $d(\tau)$ е достатъчно малък.

\subsection{Малки и големи суми на Дарбу}

Нека функцията $f$ е ограничена върху $[a,b]$ и нека
\[
\tau\colon\quad a=x_0<x_1<\cdots<x_n=b
\]
е разбиване. За всеки подинтервал полагаме
\[
m_i=\inf\{f(x)\mid x\in[x_{i-1},x_i]\},
\qquad
M_i=\sup\{f(x)\mid x\in[x_{i-1},x_i]\}.
\]
Тъй като $f$ е ограничена, числата $m_i$ и $M_i$ са крайни.

\begin{defn}
\emph{Малка сума на Дарбу} на $f$ по разбиването $\tau$ се нарича числото
\[
s_f(\tau)=\sum_{i=1}^{n}m_i(x_i-x_{i-1}).
\]
\emph{Голяма сума на Дарбу} на $f$ по разбиването $\tau$ се нарича числото
\[
S_f(\tau)=\sum_{i=1}^{n}M_i(x_i-x_{i-1}).
\]
\end{defn}

Геометрично малката сума се получава, ако над всеки подинтервал се построи правоъгълник с височина $m_i$, а голямата --- с височина $M_i$. Следователно малката сума дава долна, а голямата --- горна оценка на търсената натрупана величина.

\[
\begin{array}{c}
\text{над }[x_{i-1},x_i]\text{ височината на долния правоъгълник е }m_i,
\\[1mm]
\text{а височината на горния правоъгълник е }M_i:
\\[2mm]
\begin{array}{c}
\text{\rule{35mm}{0.4pt}}\\[-1mm]
\hspace{4mm}x_{i-1}\hspace{20mm}x_i
\end{array}
\end{array}
\]

По-точно, при добавяне на делящи точки долните правоъгълници могат да се повишат, а горните --- да се понижат.

\begin{lem}
\label{lem:refinement}
Нека $\tau^*$ се получава от разбиването $\tau$ чрез добавяне на една или повече нови делящи точки. Тогава
\[
s_f(\tau^*)\geq s_f(\tau),
\qquad
S_f(\tau^*)\leq S_f(\tau).
\]
\end{lem}

\begin{proof}
Достатъчно е да разгледаме добавянето на една нова точка, тъй като общият случай следва чрез многократно прилагане на същото разсъждение.

Нека $x^*\in(x_{i-1},x_i)$ е добавена точка. Единствено приносът на интервала $[x_{i-1},x_i]$ се променя. Нека
\[
m_i'=\inf_{x\in[x_{i-1},x^*]}f(x),
\qquad
m_i''=\inf_{x\in[x^*,x_i]}f(x).
\]
Понеже
\[
[x_{i-1},x^*]\subseteq[x_{i-1},x_i],
\qquad
[x^*,x_i]\subseteq[x_{i-1},x_i],
\]
имаме
\[
m_i'\geq m_i,
\qquad
m_i''\geq m_i.
\]
Следователно
\begin{align*}
s_f(\tau^*)-s_f(\tau)
&=m_i'(x^*-x_{i-1})+m_i''(x_i-x^*)-m_i(x_i-x_{i-1})\\
&\geq m_i(x^*-x_{i-1})+m_i(x_i-x^*)-m_i(x_i-x_{i-1})\\
&=0.
\end{align*}
Значи $s_f(\tau^*)\geq s_f(\tau)$.

За големите суми разсъждението е аналогично. Ако $M_i'$ и $M_i''$ са супремумите на функцията върху двата нови подинтервала, то
\[
M_i'\leq M_i,\qquad M_i''\leq M_i.
\]
Затова
\[
S_f(\tau^*)\leq S_f(\tau).
\]
\end{proof}

\begin{supp}[Бележка]
В част от предоставените записи означенията за горен и долен интеграл са разменени. Съгласно определенията чрез суми на Дарбу горният интеграл се задава чрез инфимум на \emph{големите} суми, а долният --- чрез супремум на \emph{малките} суми. Това е и единственото съгласувано с неравенството $s_f(\tau)\leq S_f(\tau)$ означение.
\end{supp}

\subsection{Горен и долен интеграл на Дарбу}

За всяко разбиване $\tau$ е изпълнено
\[
s_f(\tau)\leq S_f(\tau).
\]
Освен това всяка малка сума е не по-голяма от всяка голяма сума. Действително, ако $\tau_1$ и $\tau_2$ са две разбивания, разглеждаме общото им доразбиване $\tau_1\cup\tau_2$. От Лема \ref{lem:refinement} следва
\[
s_f(\tau_1)\leq s_f(\tau_1\cup\tau_2)\leq S_f(\tau_1\cup\tau_2)\leq S_f(\tau_2).
\]

\begin{defn}
\emph{Долен интеграл на Дарбу} на ограничената функция $f$ върху $[a,b]$ наричаме
\[
\underline{\int_a^b} f
=
\sup\{s_f(\tau)\mid \tau\text{ е разбиване на }[a,b]\}.
\]
\emph{Горен интеграл на Дарбу} наричаме
\[
\overline{\int_a^b} f
=
\inf\{S_f(\tau)\mid \tau\text{ е разбиване на }[a,b]\}.
\]
\end{defn}

От разгледаното неравенство между малките и големите суми следва, че
\[
\underline{\int_a^b} f\leq\overline{\int_a^b} f.
\]

\begin{defn}
Ограничената функция $f\colon[a,b]\to\mathbb{R}$ се нарича \emph{интегруема по Риман} върху $[a,b]$, ако
\[
\underline{\int_a^b} f=\overline{\int_a^b} f.
\]
Общата стойност се нарича \emph{риманов определен интеграл} на $f$ върху $[a,b]$ и се означава с
\[
\int_a^b f(x)\,dx.
\]
\end{defn}

Така определението чрез Дарбу формализира идеята, че долните и горните стъпаловидни оценки могат да бъдат направени произволно близки.

\section{Критерий за интегруемост}

\begin{thm}[Критерий на Дарбу]
\label{thm:darbourcriterion}
Нека $f\colon[a,b]\to\mathbb{R}$ е ограничена. Тогава $f$ е интегруема по Риман тогава и само тогава, когато за всяко $\varepsilon>0$ съществува разбиване $\tau$ на $[a,b]$, за което
\[
S_f(\tau)-s_f(\tau)<\varepsilon.
\]
\end{thm}

\begin{proof}
Нека първо $f$ е интегруема и
\[
I=\int_a^b f(x)\,dx
=
\underline{\int_a^b}f
=
\overline{\int_a^b}f.
\]
Нека $\varepsilon>0$. От свойството на супремума съществува разбиване $\tau_1$, за което
\[
s_f(\tau_1)>I-\frac{\varepsilon}{2}.
\]
От свойството на инфимума съществува разбиване $\tau_2$, за което
\[
S_f(\tau_2)<I+\frac{\varepsilon}{2}.
\]
Нека $\tau=\tau_1\cup\tau_2$ е общото доразбиване. Според Лема \ref{lem:refinement},
\[
s_f(\tau)\geq s_f(\tau_1),
\qquad
S_f(\tau)\leq S_f(\tau_2).
\]
Следователно
\[
S_f(\tau)-s_f(\tau)
\leq S_f(\tau_2)-s_f(\tau_1)
<
\left(I+\frac{\varepsilon}{2}\right)
-
\left(I-\frac{\varepsilon}{2}\right)
=
\varepsilon.
\]

Обратно, нека за всяко $\varepsilon>0$ има разбиване $\tau$, за което
\[
S_f(\tau)-s_f(\tau)<\varepsilon.
\]
За всяко разбиване е изпълнено
\[
s_f(\tau)\leq\underline{\int_a^b}f
\leq\overline{\int_a^b}f
\leq S_f(\tau).
\]
Затова
\[
0\leq
\overline{\int_a^b}f-\underline{\int_a^b}f
\leq S_f(\tau)-s_f(\tau)
<\varepsilon.
\]
Това е вярно за всяко $\varepsilon>0$, откъдето
\[
\overline{\int_a^b}f=\underline{\int_a^b}f.
\]
Следователно $f$ е интегруема по Риман.
\end{proof}

\section{Интегруемост на непрекъснатите функции}

Ще използваме без доказателство теоремата на Кантор: всяка непрекъсната функция върху краен затворен интервал е равномерно непрекъсната. Ще използваме също, че непрекъсната функция върху $[a,b]$ е ограничена.

\begin{thm}
Всяка непрекъсната функция $f\colon[a,b]\to\mathbb{R}$ е интегруема по Риман.
\end{thm}

\begin{proof}
Нека $\varepsilon>0$. Понеже $f$ е равномерно непрекъсната върху $[a,b]$, съществува $\delta>0$, такова че за всички $x',x''\in[a,b]$ от
\[
|x'-x''|<\delta
\]
следва
\[
|f(x')-f(x'')|<\frac{\varepsilon}{b-a}.
\]

Избираме разбиване
\[
\tau\colon\quad a=x_0<x_1<\cdots<x_n=b,
\qquad
d(\tau)<\delta.
\]
За произволен индекс $i$ и произволни точки
\[
x',x''\in[x_{i-1},x_i]
\]
имаме
\[
|x'-x''|\leq x_i-x_{i-1}\leq d(\tau)<\delta.
\]
Следователно
\[
|f(x')-f(x'')|<\frac{\varepsilon}{b-a}.
\]
Оттук осцилацията на $f$ върху всеки подинтервал удовлетворява
\[
M_i-m_i\leq\frac{\varepsilon}{b-a}.
\]
Тогава
\begin{align*}
S_f(\tau)-s_f(\tau)
&=\sum_{i=1}^{n}(M_i-m_i)(x_i-x_{i-1})\\
&\leq
\frac{\varepsilon}{b-a}
\sum_{i=1}^{n}(x_i-x_{i-1})\\
&=
\frac{\varepsilon}{b-a}(b-a)
=
\varepsilon.
\end{align*}
По Теорема \ref{thm:darbourcriterion} функцията $f$ е интегруема по Риман.
\end{proof}

\section{Основни свойства на римановия интеграл}

Следващите свойства се използват без доказателство. Нека разглежданите функции са интегруеми по Риман върху съответните интервали.

\begin{prop}[Линейност]
Ако $f,g\colon[a,b]\to\mathbb{R}$ са интегруеми и $\lambda,\mu\in\mathbb{R}$, то $\lambda f+\mu g$ е интегруема и
\[
\int_a^b\bigl(\lambda f(x)+\mu g(x)\bigr)\,dx
=
\lambda\int_a^b f(x)\,dx
+
\mu\int_a^b g(x)\,dx.
\]
В частност,
\[
\int_a^b(f(x)+g(x))\,dx
=
\int_a^b f(x)\,dx+\int_a^b g(x)\,dx,
\]
\[
\int_a^b\lambda f(x)\,dx
=
\lambda\int_a^b f(x)\,dx.
\]
\end{prop}

\begin{prop}[Адитивност по интервала]
Ако $c\in(a,b)$, то $f$ е интегруема върху $[a,b]$ тогава и само тогава, когато е интегруема върху $[a,c]$ и $[c,b]$. В този случай
\[
\int_a^b f(x)\,dx
=
\int_a^c f(x)\,dx+\int_c^b f(x)\,dx.
\]
\end{prop}

\begin{prop}[Смяна на границите]
За интегруема функция $f$ е изпълнено
\[
\int_b^a f(x)\,dx=-\int_a^b f(x)\,dx.
\]
\end{prop}

\begin{prop}[Позитивност и интегриране на неравенства]
Ако $f(x)\geq0$ за всяко $x\in[a,b]$, то
\[
\int_a^b f(x)\,dx\geq0.
\]
По-общо, ако $f(x)\leq g(x)$ за всяко $x\in[a,b]$, то
\[
\int_a^b f(x)\,dx\leq\int_a^b g(x)\,dx.
\]
\end{prop}

\begin{cor}
Ако $f$ е интегруема върху $[a,b]$, то $|f|$ също е интегруема и
\[
\left|\int_a^b f(x)\,dx\right|
\leq
\int_a^b |f(x)|\,dx.
\]
\end{cor}

\begin{example}
За константната функция $f(x)=k$ е изпълнено
\[
\int_a^b k\,dx=k(b-a).
\]
\end{example}

\section{Теорема за средната стойност за интеграли}

\begin{thm}[Теорема за средната стойност]
\label{thm:meanintegral}
Нека $f\colon[a,b]\to\mathbb{R}$ е непрекъсната. Тогава съществува $c\in[a,b]$, такова че
\[
\int_a^b f(x)\,dx=f(c)(b-a).
\]
\end{thm}

\begin{proof}
Понеже $f$ е непрекъсната върху затворения интервал $[a,b]$, тя приема най-малка и най-голяма стойност. Нека
\[
m=\min_{x\in[a,b]}f(x),
\qquad
M=\max_{x\in[a,b]}f(x).
\]
Тогава
\[
m\leq f(x)\leq M,
\qquad x\in[a,b].
\]
Чрез свойството за интегриране на неравенства получаваме
\[
\int_a^b m\,dx
\leq
\int_a^b f(x)\,dx
\leq
\int_a^b M\,dx.
\]
Следователно
\[
m(b-a)
\leq
\int_a^b f(x)\,dx
\leq
M(b-a).
\]
Тъй като $b-a>0$,
\[
m\leq
\frac{1}{b-a}\int_a^b f(x)\,dx
\leq M.
\]
Непрекъснатата функция приема всички стойности между своя минимум и максимум. Следователно съществува $c\in[a,b]$, за което
\[
f(c)=\frac{1}{b-a}\int_a^b f(x)\,dx.
\]
Умножаваме по $b-a$ и получаваме
\[
\int_a^b f(x)\,dx=f(c)(b-a).
\]
\end{proof}

Числото
\[
\frac{1}{b-a}\int_a^b f(x)\,dx
\]
се нарича средна стойност на функцията $f$ върху $[a,b]$. Теоремата показва, че за непрекъсната функция тази средна стойност действително се достига като стойност $f(c)$ в поне една точка от интервала.

\section{Теорема на Нютон--Лайбниц}

\begin{keythm}[Теорема на Нютон--Лайбниц]
\label{thm:newtonleibniz}
Нека $f\colon[a,b]\to\mathbb{R}$ е непрекъсната и
\[
F(x)=\int_a^x f(t)\,dt,
\qquad x\in[a,b].
\]
Тогава $F$ е диференцируема във всяка вътрешна точка на $[a,b]$ и
\[
F'(x)=f(x).
\]
В краищата $a$ и $b$ същото равенство се разбира чрез съответната едностранна производна.
\end{keythm}

\begin{proof}
Нека $x_0\in[a,b]$ е фиксирана точка и $h\ne0$ е такова, че
\[
x_0+h\in[a,b].
\]
От адитивността на интеграла следва
\[
F(x_0+h)-F(x_0)
=
\int_{x_0}^{x_0+h}f(t)\,dt.
\]
Следователно
\[
\frac{F(x_0+h)-F(x_0)}{h}
=
\frac{1}{h}\int_{x_0}^{x_0+h}f(t)\,dt.
\]
Изваждаме $f(x_0)$ и използваме, че
\[
f(x_0)=\frac{1}{h}\int_{x_0}^{x_0+h}f(x_0)\,dt.
\]
Получаваме
\[
\frac{F(x_0+h)-F(x_0)}{h}-f(x_0)
=
\frac{1}{h}\int_{x_0}^{x_0+h}
\bigl(f(t)-f(x_0)\bigr)\,dt.
\]

Нека $\varepsilon>0$. От непрекъснатостта на $f$ в $x_0$ съществува $\delta>0$, такова че
\[
|t-x_0|<\delta
\quad\Longrightarrow\quad
|f(t)-f(x_0)|<\varepsilon.
\]
Ако $|h|<\delta$, то всяка точка $t$ между $x_0$ и $x_0+h$ удовлетворява
\[
|t-x_0|\leq |h|<\delta.
\]

При $h>0$, използвайки неравенството за интеграла, имаме
\begin{align*}
\left|
\frac{F(x_0+h)-F(x_0)}{h}-f(x_0)
\right|
&\leq
\frac{1}{h}
\int_{x_0}^{x_0+h}|f(t)-f(x_0)|\,dt\\
&\leq
\frac{1}{h}\int_{x_0}^{x_0+h}\varepsilon\,dt
=
\varepsilon.
\end{align*}
При $h<0$ аналогично,
\begin{align*}
\left|
\frac{F(x_0+h)-F(x_0)}{h}-f(x_0)
\right|
&\leq
\frac{1}{|h|}
\int_{x_0+h}^{x_0}|f(t)-f(x_0)|\,dt\\
&\leq
\frac{1}{|h|}
\int_{x_0+h}^{x_0}\varepsilon\,dt
=
\varepsilon.
\end{align*}

Следователно
\[
\lim_{h\to0}
\frac{F(x_0+h)-F(x_0)}{h}
=
f(x_0).
\]
Значи $F'(x_0)=f(x_0)$. Тъй като $x_0$ е произволна точка, твърдението е доказано.
\end{proof}

Теоремата показва, че функцията
\[
x\longmapsto\int_a^x f(t)\,dt
\]
е примитивна на непрекъснатата функция $f$.

\subsection{Изчисляване на определени интеграли}

\begin{cor}[Формула на Нютон--Лайбниц]
Нека $f\colon[a,b]\to\mathbb{R}$ е непрекъсната и $\Phi$ е примитивна на $f$ върху $[a,b]$, тоест
\[
\Phi'(x)=f(x).
\]
Тогава
\[
\int_a^b f(x)\,dx=\Phi(b)-\Phi(a).
\]
Обичайното означение е
\[
\int_a^b f(x)\,dx=\left.\Phi(x)\right|_a^b.
\]
\end{cor}

\begin{proof}
По Теорема \ref{thm:newtonleibniz} функцията
\[
F(x)=\int_a^x f(t)\,dt
\]
също е примитивна на $f$. Следователно за някоя константа $C$ е изпълнено
\[
\Phi(x)=F(x)+C.
\]
При $x=a$ получаваме
\[
\Phi(a)=F(a)+C=C,
\]
защото
\[
F(a)=\int_a^a f(t)\,dt=0.
\]
Следователно
\[
\Phi(x)=\int_a^x f(t)\,dt+\Phi(a),
\]
или
\[
\int_a^x f(t)\,dt=\Phi(x)-\Phi(a).
\]
При $x=b$ получаваме търсената формула.
\end{proof}

\begin{example}
Да се изчисли
\[
\int_0^1 3x^2\,dx.
\]
Примитивна на $3x^2$ е $\Phi(x)=x^3$. Следователно
\[
\int_0^1 3x^2\,dx
=
\left.x^3\right|_0^1
=
1^3-0^3
=
1.
\]
\end{example}

Практическият алгоритъм за пресмятане на определен интеграл на непрекъсната функция е следният:

\begin{enumerate}
\item Намира се примитивна функция $\Phi$ на подинтегралната функция $f$.
\item Изчисляват се стойностите $\Phi(a)$ и $\Phi(b)$.
\item Прилага се формулата
\[
\int_a^b f(x)\,dx=\Phi(b)-\Phi(a).
\]
\end{enumerate}

\section{Изпитен фокус}

\begin{itemize}
\item Да се дефинират разбиване на $[a,b]$, делящи точки и диаметър $d(\tau)$.
\item Да се дефинират $m_i$, $M_i$, малка сума $s_f(\tau)$ и голяма сума $S_f(\tau)$ на Дарбу.
\item Да се възпроизведе доказателството, че при доразбиване малките суми не намаляват, а големите не нарастват.
\item Да се дефинират долен и горен интеграл на Дарбу и интегруемост по Риман чрез равенството им.
\item Да се докаже критерият на Дарбу:
\[
f\text{ е интегруема}
\Longleftrightarrow
\forall\varepsilon>0\ \exists\tau:
S_f(\tau)-s_f(\tau)<\varepsilon.
\]
\item Да се обясни доказателството за интегруемост на непрекъсната функция: теорема на Кантор, избор на разбиване с малък диаметър и оценка на $S_f(\tau)-s_f(\tau)$.
\item Да се изброят линейност, адитивност, позитивност, интегриране на неравенства, смяна на границите и неравенството с $|f|$.
\item Да се докаже теоремата за средната стойност за интеграли чрез минимум, максимум, интегриране на неравенства и свойството за междинните стойности.
\item Да се докаже теоремата на Нютон--Лайбниц чрез разликата
\[
\frac{1}{h}\int_x^{x+h}\bigl(f(t)-f(x)\bigr)\,dt
\]
и непрекъснатостта на $f$.
\item Да се използва формулата
\[
\int_a^b f(x)\,dx=\Phi(b)-\Phi(a)
\]
за изчисляване на определен интеграл чрез примитивна функция.
\end{itemize}
